\documentclass[11pt, a4paper]{article}

% ── Codificación y Lenguaje ──────────────────────────────────────────────────
\usepackage[utf8]{inputenc}
\usepackage[T1]{fontenc}
\usepackage[spanish, es-tabla]{babel}

% ── Geometría ───────────────────────────────────────────────────────────────
\usepackage[top=2.2cm, bottom=2.2cm, left=2.2cm, right=2.2cm]{geometry}

% ── Matemáticas ─────────────────────────────────────────────────────────────
\usepackage{amsmath, amssymb, amsthm}
\usepackage{mathtools}

% ── Tablas ──────────────────────────────────────────────────────────────────
\usepackage{booktabs}
\usepackage{array}
\usepackage{multirow}

% ── Listas ──────────────────────────────────────────────────────────────────
\usepackage{enumitem}

% ── Encabezado y Pie de Página ──────────────────────────────────────────────
\usepackage{fancyhdr}

% ── Gráficos y Diagramas ────────────────────────────────────────────────────
\usepackage{tikz}
\usetikzlibrary{shapes.geometric, arrows.meta, positioning, calc}
\usepackage{graphicx}

% ── Espaciado ───────────────────────────────────────────────────────────────
\usepackage{setspace}
\setlength{\parskip}{5pt}
\setlength{\parindent}{0pt}

% Encabezado y Pie de Página
\pagestyle{fancy}
\fancyhf{}
\fancyhead[L]{\small\textbf{ICS3213 Gestión de Operaciones} --- 1er Semestre 2026}
\fancyhead[R]{\small Interrogación 2 (Desarrollo)}
\fancyfoot[C]{\small\thepage}
\renewcommand{\headrulewidth}{0.4pt}

\begin{document}

\begin{center}
  \textbf{\Large Pontificia Universidad Católica de Chile} \\
  \textbf{Escuela de Ingeniería --- Departamento de Ingeniería Industrial y de Sistemas} \\
  \vspace{0.4cm}
  {\Large\bfseries Interrogación 2 (Para la Casa) --- Respuestas de Desarrollo} \\
  \vspace{0.3cm}
  \begin{tabular}{rl}
    \textbf{Nombre:} & Fabián Ignacio Ortega Llantén \\
    \textbf{Email UC:} & \texttt{fgortega@uc.cl} \\
    \textbf{Sección:} & 1
  \end{tabular}
\end{center}

\vspace{0.3cm}
\rule{\textwidth}{0.4pt}
\vspace{0.5cm}

% ════════════════════════════════════════════════════════════════════════════
% PREGUNTA 1
% ════════════════════════════════════════════════════════════════════════════
\section*{Pregunta 1: Planificación Agregada y MRP}

\subsection*{i. Plan Agregado y Cantidad de Personal (Estrategia de Nivelación)}

Para diseñar un plan de nivelación con producción mensual constante $P$ y fuerza laboral constante $W$ sin permitir faltantes ($I_t \ge 0$), planteamos las restricciones de balance de inventario acumulado para el trimestre (octubre-diciembre), partiendo de un inventario inicial $I_0 = 0$:
\begin{align*}
  t = 1 \quad (\text{Oct}): \quad & I_1 = I_0 + P - D_1 = P - 800 \ge 0 \implies P \ge 800 \\
  t = 2 \quad (\text{Nov}): \quad & I_2 = I_1 + P - D_2 = 2P - 2.000 \ge 0 \implies P \ge 1.000 \\
  t = 3 \quad (\text{Dic}): \quad & I_3 = I_2 + P - D_3 = 3P - 3.600 \ge 0 \implies P \ge 1.200
\end{align*}

Para evitar quiebres de stock en todos los meses, la producción constante mínima debe ser:
\[
  P^* = \max(800, 1.000, 1.200) = \mathbf{1.200 \text{ unidades/mes}}
\]

La fuerza laboral constante ($W$) necesaria se calcula dividiendo la producción mensual entre la productividad de un trabajador ($10 \text{ drones/semana} \times 4 \text{ semanas/mes} = 40 \text{ drones/trabajador-mes}$):
\[
  W = \frac{1.200 \text{ unidades/mes}}{40 \text{ unidades/trabajador-mes}} = \mathbf{30 \text{ trabajadores}}
\]

El balance de inventarios y sus costos asociados al final de cada mes (con $H = \$50$ dólares/dron-mes) son:
\begin{itemize}
  \item \textbf{Octubre:} $I_1 = 0 + 1.200 - 800 = 400 \text{ unidades} \implies \text{Costo} = 400 \times \$50 = \$20.000$.
  \item \textbf{Noviembre:} $I_2 = 400 + 1.200 - 1.200 = 400 \text{ unidades} \implies \text{Costo} = 400 \times \$50 = \$20.000$.
  \item \textbf{Diciembre:} $I_3 = 400 + 1.200 - 1.600 = 0 \text{ unidades} \implies \text{Costo} = 0 \times \$50 = \$0$.
\end{itemize}

El costo total de mantener inventario para el trimestre es:
\[
  \text{Costo Total} = \$20.000 + \$20.000 + \$0 = \mathbf{\$40.000}
\]

\subsection*{ii. Planificación MRP para el Componente MCM (Octubre)}

La producción del modelo estrella corresponde al 50\% del total agregado de octubre ($1.200 \times 0.5 = 600$ drones estrella/mes). Se distribuye de manera uniforme durante las 4 semanas de octubre: $600 / 4 = 150$ drones estrella/semana en el Programa Maestro de Producción (MPS). Dado que cada dron requiere 1 MCM, los requerimientos brutos ($GR$) de MCM son de $150$ unidades/semana.

A partir de un inventario inicial de $50$ unidades de MCM, una recepción programada de $200$ en la Semana 1, un Lead Time ($LT$) de 2 semanas y una política de loteo Lote a Lote (L4L), se construye la siguiente matriz MRP:

\begin{center}
  \small
  \begin{tabular}{lccccc}
    \toprule
    \textbf{Concepto / Semana} & \textbf{Semana 0} & \textbf{Semana 1} & \textbf{Semana 2} & \textbf{Semana 3} & \textbf{Semana 4} \\
    \midrule
    Requerimiento Bruto ($GR$) & & 150 & 150 & 150 & 150 \\
    Recepciones Programadas ($SR$) & & 200 & 0 & 0 & 0 \\
    Stock Disponible Proyectado ($I$) & 50 & 100 & 0 & 0 & 0 \\
    Requerimientos Netos ($NR$) & & 0 & 50 & 150 & 150 \\
    Recepciones Planificadas ($POR$) & & 0 & 50 & 150 & 150 \\
    Lanzamientos Planificados ($PORelease$) & 50 & 150 & 150 & 0 & 0 \\
    \bottomrule
  \end{tabular}
\end{center}

\textbf{Lanzamientos requeridos durante el mes de octubre (de la Semana 1 a la Semana 4):}
\begin{itemize}
  \item En la \textbf{Semana 1}, se deben emitir pedidos por \textbf{150 unidades} de MCM (que llegarán en la Semana 3).
  \item En la \textbf{Semana 2}, se deben emitir pedidos por \textbf{150 unidades} de MCM (que llegarán en la Semana 4).
\end{itemize}
\textit{Nota: El lanzamiento de 50 unidades de la Semana 0 se debió emitir antes de comenzar el mes de octubre.}

\subsection*{iii. Análisis Crítico EOQ vs. MRP}

El modelo **EOQ clásico** es inadecuado para la planificación de materiales en un sistema MRP por las siguientes razones matemáticas y operativas:
\begin{enumerate}
  \item **Demanda Discontinua e Intermitente:** EOQ asume que la demanda es continua, conocida y constante a una tasa uniforme. En cambio, los requerimientos de componentes de un MRP son de demanda dependiente, la cual es discreta, altamente variable e intermitente (*lumpy demand*) debido a la consolidación de lotes en niveles superiores.
  \item **Costos Innecesarios de Inventario:** Si se fuerza un lote fijo ($Q^*_{\text{EOQ}}$) en períodos de baja demanda, se generará un inventario residual innecesario que incrementa notablemente los costos de mantener stock.
  \item **Riesgo de Desabastecimiento:** En semanas con requerimientos concentrados altos que superen la cantidad fija de EOQ, el lote rígido puede no cubrir la necesidad neta, provocando retrasos en la producción.
\end{enumerate}

**Recomendación:** Se aconseja utilizar el algoritmo de **Wagner-Whitin (WW)** o la heurística de **Silver-Meal (SM)**. Estos métodos son dinámicos y buscan balancear los costos de pedir y mantener adaptándose a la variabilidad de la demanda en cada período.

\newpage

% ════════════════════════════════════════════════════════════════════════════
% PREGUNTA 2
% ════════════════════════════════════════════════════════════════════════════
\section*{Pregunta 2: Administración de Proyectos PERT/CPM}

\subsection*{i. Diagrama PERT, Ruta Crítica y Estadísticas del Proyecto}

Utilizando los datos de tiempos esperados de las actividades, determinamos los tiempos tempranos ($ES, EF$) y tardíos ($LS, LF$) mediante las pasadas hacia adelante y hacia atrás de la metodología CPM. La holgura se calcula como $H = LS - ES = LF - EF$.

La siguiente tabla resume los cálculos detallados para cada actividad:
\begin{center}
  \small
  \begin{tabular}{ccccccccc}
    \toprule
    Actividad & Predecesor & Duración ($TE$) & $ES$ & $EF$ & $LS$ & $LF$ & Holgura ($H$) & Crítica \\
    \midrule
    A & - & 3 & 0 & 3 & 0 & 3 & 0 & Sí \\
    B & A & 5 & 3 & 8 & 3 & 8 & 0 & Sí \\
    C & A & 4 & 3 & 7 & 5 & 9 & 2 & No \\
    D & B & 6 & 8 & 14 & 8 & 14 & 0 & Sí \\
    E & B, C & 5 & 8 & 13 & 9 & 14 & 1 & No \\
    F & C & 7 & 7 & 14 & 11 & 18 & 4 & No \\
    G & D, E & 4 & 14 & 18 & 14 & 18 & 0 & Sí \\
    H & F, G & 3 & 18 & 21 & 18 & 21 & 0 & Sí \\
    \bottomrule
  \end{tabular}
\end{center}

Las actividades con holgura cero definen la **Ruta Crítica**, la cual corresponde a **A - B - D - G - H**.

A continuación se presenta el diagrama del proyecto en formato de Actividad en Nodo (AON) indicando el camino crítico (en rojo y línea gruesa) y las correspondientes pasadas CPM:

\begin{center}
\begin{tikzpicture}[
  node distance=1.5cm and 2cm,
  activity/.style={
    draw, rectangle, rounded corners=2pt,
    fill=gray!10, text width=2.5cm, align=center,
    font=\small, minimum height=1.3cm
  },
  critical/.style={
    draw=red!80!black, thick, rectangle, rounded corners=2pt,
    fill=red!5, text width=2.5cm, align=center,
    font=\small, minimum height=1.3cm
  },
  arrow/.style={-Stealth, thick, draw=gray!80},
  critarrow/.style={-Stealth, ultra thick, draw=red!80!black}
]

  % Nodes
  \node[critical] (A) at (0, 0) {
    \textbf{A} ($TE=3$) \\
    \footnotesize ES: 0 \hfill EF: 3 \\
    \footnotesize LS: 0 \hfill LF: 3
  };
  
  \node[critical] (B) at (3.8, 1.8) {
    \textbf{B} ($TE=5$) \\
    \footnotesize ES: 3 \hfill EF: 8 \\
    \footnotesize LS: 3 \hfill LF: 8
  };
  
  \node[activity] (C) at (3.8, -1.8) {
    \textbf{C} ($TE=4$) \\
    \footnotesize ES: 3 \hfill EF: 7 \\
    \footnotesize LS: 5 \hfill LF: 9
  };
  
  \node[critical] (D) at (7.6, 1.8) {
    \textbf{D} ($TE=6$) \\
    \footnotesize ES: 8 \hfill EF: 14 \\
    \footnotesize LS: 8 \hfill LF: 14
  };
  
  \node[activity] (E) at (7.6, 0) {
    \textbf{E} ($TE=5$) \\
    \footnotesize ES: 8 \hfill EF: 13 \\
    \footnotesize LS: 9 \hfill LF: 14
  };
  
  \node[activity] (F) at (7.6, -1.8) {
    \textbf{F} ($TE=7$) \\
    \footnotesize ES: 7 \hfill EF: 14 \\
    \footnotesize LS: 11 \hfill LF: 18
  };
  
  \node[critical] (G) at (11.4, 0.9) {
    \textbf{G} ($TE=4$) \\
    \footnotesize ES: 14 \hfill EF: 18 \\
    \footnotesize LS: 14 \hfill LF: 18
  };
  
  \node[critical] (H) at (14.8, 0) {
    \textbf{H} ($TE=3$) \\
    \footnotesize ES: 18 \hfill EF: 21 \\
    \footnotesize LS: 18 \hfill LF: 21
  };

  % Connections
  \draw[critarrow] (A) -- (B);
  \draw[arrow] (A) -- (C);
  \draw[critarrow] (B) -- (D);
  \draw[arrow] (B) -- (E);
  \draw[arrow] (C) -- (E);
  \draw[arrow] (C) -- (F);
  \draw[critarrow] (D) -- (G);
  \draw[arrow] (E) -- (G);
  \draw[arrow] (F) -- (H);
  \draw[critarrow] (G) -- (H);

\end{tikzpicture}
\end{center}

* **Duración esperada del proyecto ($E[T]$):**
  \[
    E[T] = TE_A + TE_B + TE_D + TE_G + TE_H = 3 + 5 + 6 + 4 + 3 = \mathbf{21 \text{ semanas}}
  \]
* **Varianza del proyecto ($\text{Var}[T]$):**
  \[
    \text{Var}[T] = \sigma_A^2 + \sigma_B^2 + \sigma_D^2 + \sigma_G^2 + \sigma_H^2 = \frac{4}{9} + \frac{4}{9} + \frac{16}{9} + \frac{4}{9} + \frac{1}{9} = \frac{29}{9} \approx \mathbf{3.222 \text{ semanas}^2}
  \]
* **Desviación estándar del proyecto ($\sigma_p$):**
  \[
    \sigma_p = \sqrt{3.222} \approx \mathbf{1.795 \text{ semanas}}
  \]

\subsection*{ii. Probabilidad e Intervalo de Confianza}

* **Probabilidad de terminar en o antes de 23 semanas ($P(T \le 23)$):**
  Aproximando la duración del proyecto a una distribución normal por el TLC, normalizamos para $T = 23$:
  \[
    Z = \frac{23 - E[T]}{\sigma_p} = \frac{23 - 21}{1.795} = 1.114
  \]
  Buscando en la tabla de distribución normal estándar para $Z = 1.11$, obtenemos $\Phi(1.11) \approx 0.8665$.
  Utilizando interpolación lineal para un valor más preciso de $Z = 1.114$:
  \[
    P(T \le 23) \approx \Phi(1.11) + 0.4 \times (\Phi(1.12) - \Phi(1.11)) \approx 0.8665 + 0.4 \times (0.8686 - 0.8665) = \mathbf{86.73\%}
  \]

* **Intervalo de confianza al 95\% (dos colas):**
  El valor de $Z$ crítico para el 95\% de confianza es $Z_{\alpha/2} = 1.96$:
  \[
    IC = E[T] \pm 1.96 \cdot \sigma_p = 21 \pm 1.96 \cdot (1.795) = 21 \pm 3.518 \implies \mathbf{[17.48, 24.52] \text{ semanas}}
  \]

\subsection*{iii. Evaluación de Contrato (Bono y Penalización Fija)}

Evaluamos financieramente la conveniencia de aceptar el contrato mediante el cálculo de su Valor Esperado ($VE$):
\begin{itemize}
  \item Bono: $\$1.5$ millones si finaliza en o antes de 21 semanas ($T \le 21$).
    \[
      Z_{\text{bono}} = \frac{21 - 21}{\sigma_p} = 0 \implies P(T \le 21) = 0.50
    \]
  \item Penalización: $\$4.0$ millones si finaliza después de 23 semanas ($T > 23$).
    \[
      P(T > 23) = 1 - P(T \le 23) = 1 - 0.8673 = 0.1327
    \]
\end{itemize}

El Valor Esperado de la utilidad del contrato es:
\[
  VE = \$1.5 \cdot P(T \le 21) - \$4.0 \cdot P(T > 23) = 1.5(0.50) - 4.0(0.1327) = 0.75 - 0.5308 = \mathbf{+0.2192 \text{ millones}}
\]

**Conclusión:** Dado que el Valor Esperado es positivo ($+\$219.200$), **se debe aceptar el contrato**.

\subsection*{iv. Modelo Matemático para Bono/Penalización Lineal}

Si se introduce un bono lineal de $\$1.5$ millones por semana de adelanto ($T \le 21$) y una penalización lineal de $\$4.0$ millones por semana de atraso ($T > 23$), el resultado financiero del contrato $f(T)$ es una función continua por tramos de la variable aleatoria de finalización $T \sim N(21, 3.222)$:
\[
  f(T) = \begin{cases}
    1.5 \cdot (21 - T) & \text{si } T \le 21 \\
    0 & \text{si } 21 < T \le 23 \\
    -4.0 \cdot (T - 23) & \text{si } T > 23
  \end{cases}
\]

La condición matemática para aceptar el contrato es que su valor esperado sea estrictamente positivo ($E[f(T)] > 0$). Denotando la función de densidad de probabilidad de la duración del proyecto como $\phi_T(t)$, la condición se plantea analíticamente como:
\[
  \int_{-\infty}^{21} 1.5(21 - t) \phi_T(t) dt - \int_{23}^{\infty} 4.0(t - 23) \phi_T(t) dt > 0
\]
Reemplazando la función de densidad normal $\phi_T(t) = \frac{1}{\sqrt{2\pi}\sigma_p} e^{-\frac{(t-E[T])^2}{2\sigma_p^2}}$:
\[
  \int_{-\infty}^{21} 1.5(21 - t) \frac{1}{\sqrt{2\pi}(1.795)} e^{-\frac{(t-21)^2}{2(3.222)}} dt - \int_{23}^{\infty} 4.0(t - 23) \frac{1}{\sqrt{2\pi}(1.795)} e^{-\frac{(t-21)^2}{2(3.222)}} dt > 0
\]

\newpage

% ════════════════════════════════════════════════════════════════════════════
% PREGUNTA 3
% ════════════════════════════════════════════════════════════════════════════
\section*{Pregunta 3: Gestión de Capacidad y Variabilidad}

\subsection*{a) Utilización, Tiempo de Espera en Cola y Sala de Espera}

\textbf{Datos iniciales:}
\begin{itemize}
  \item Tasa de llegadas ($\lambda$): $20$ personas/hora $\implies \lambda = \frac{20}{60} = \frac{1}{3}$ personas/minuto.
  \item Tiempo promedio de servicio ($1/\mu$): $2.4$ minutos $\implies \mu = \frac{1}{2.4} = \frac{25}{60}$ personas/minuto ($25$ personas/hora).
  \item Coeficiente de variación de arribos ($C_a$): $0.5$.
  \item Coeficiente de variación del servicio ($C_s$): $1.5$.
\end{itemize}

El nivel de utilización de la oficina ($\rho$) se define como:
\[
  \rho = \frac{\lambda}{\mu} = \frac{20}{25} = \mathbf{0.80 \quad (80\%)}
\]

Para calcular el tiempo medio de espera en cola ($W_q$), aplicamos la aproximación de Kingman para colas generales $G/G/1$:
\[
  W_q = \left( \frac{C_a^2 + C_s^2}{2} \right) \left( \frac{\rho}{1 - \rho} \right) \left( \frac{1}{\mu} \right)
\]
Reemplazando los valores correspondientes:
\[
  W_q = \left( \frac{0.5^2 + 1.5^2}{2} \right) \left( \frac{0.80}{1 - 0.80} \right) (2.4) = \left( \frac{0.25 + 2.25}{2} \right) (4) (2.4) = (1.25) \cdot 4 \cdot 2.4 = \mathbf{12 \text{ minutos}}
\]

La cantidad promedio de clientes en la cola ($L_q$) se obtiene aplicando la Ley de Little:
\[
  L_q = \lambda \cdot W_q = 20 \text{ personas/hora} \times \frac{12 \text{ minutos}}{60 \text{ minutos/hora}} = 20 \times 0.2 = \mathbf{4 \text{ personas}}
\]
Por lo tanto, la sala de espera requiere en promedio **4 asientos**.

\subsection*{b) Capacidad de Agendamiento en Agenda Digital}

Al implementar una agenda digital en bloques de 1 hora, los clientes llegan de forma determinística en los instantes agendados, eliminando por completo la variabilidad de llegada. Por lo tanto, el coeficiente de variación de los arribos es $C_a = 0$.

Buscamos la tasa de llegada máxima $\lambda$ (y la utilización asociada $\rho$) que asegure un tiempo de espera en cola menor o igual a 6 minutos ($W_q \le 6$):
\[
  W_q = \left( \frac{0^2 + 1.5^2}{2} \right) \left( \frac{\rho}{1 - \rho} \right) \left( 2.4 \right) \le 6
\]
\[
  1.125 \cdot 2.4 \cdot \left( \frac{\rho}{1 - \rho} \right) \le 6 \implies 2.7 \left( \frac{\rho}{1 - \rho} \right) \le 6
\]
Despejando la relación de utilización:
\[
  \frac{\rho}{1 - \rho} \le \frac{6}{2.7} = \frac{20}{9} \implies 9\rho \le 20(1-\rho) \implies 29\rho \le 20 \implies \rho \le \frac{20}{29} \approx 0.6897
\]

Calculamos la tasa de agendamientos máxima $\lambda$ por hora a partir de $\rho \le \frac{20}{29}$ y $\mu = 25$ clientes/hora:
\[
  \lambda = \rho \cdot \mu \le \frac{20}{29} \cdot 25 \approx 17.24 \text{ clientes/hora}
\]
Por lo tanto, para garantizar que la espera sea menor a 6 minutos, se pueden agendar a lo sumo **17 clientes por hora**.

\subsection*{c) Evaluación de Opciones de Mejora}

En la situación base, el tiempo medio de espera es $W_{q,\text{inicial}} = 12 \text{ minutos}$. Se evalúa el beneficio neto de cada opción de mejora (reducción de espera total multiplicada por $\$1.000/\text{min}$ menos el costo del proyecto):

\begin{itemize}
  \item \textbf{Opción i: Capacitación al Personal}
    \begin{itemize}
      \item Parámetros: $\rho = 0.80$, $C_a = 0.5$, nuevo $C_s = 1.0$. Costo = $\$4.000$.
      \item Nuevo tiempo de espera en cola ($W_{q,\text{Cap}}$):
        \[
          W_{q,\text{Cap}} = \left( \frac{0.5^2 + 1.0^2}{2} \right) \left( \frac{0.80}{1 - 0.80} \right) (2.4) = (0.625) \cdot 4 \cdot 2.4 = \mathbf{6 \text{ minutos}}
        \]
      \item Reducción de tiempo ($\Delta W_q$): $12 - 6 = 6$ minutos.
      \item Beneficio neto: $6 \text{ min} \times \$1.000/\text{min} - \$4.000 = \$6.000 - \$4.000 = \mathbf{+\$2.000}$.
    \end{itemize}

  \item \textbf{Opción ii: Digitalización de Servicios}
    \begin{itemize}
      \item Parámetros: nuevo $\lambda = 15 \text{ clientes/hora} \implies \rho = \frac{15}{25} = 0.60$, $C_a = 0.5$, $C_s = 1.5$. Costo = $\$6.000$.
      \item Nuevo tiempo de espera en cola ($W_{q,\text{Dig}}$):
        \[
          W_{q,\text{Dig}} = \left( \frac{0.5^2 + 1.5^2}{2} \right) \left( \frac{0.60}{1 - 0.60} \right) (2.4) = (1.25) \cdot 1.5 \cdot 2.4 = \mathbf{4.5 \text{ minutos}}
        \]
      \item Reducción de tiempo ($\Delta W_q$): $12 - 4.5 = 7.5$ minutos.
      \item Beneficio neto: $7.5 \text{ min} \times \$1.000/\text{min} - \$6.000 = \$7.500 - \$6.000 = \mathbf{+\$1.500}$.
    \end{itemize}
\end{itemize}

**Conclusión:** Se prefiere la **Opción i** (Capacitación al personal), ya que genera un beneficio neto de **$+\$2.000$**, el cual es superior al de la Opción ii ($+\$1.500$).

\end{document}
